\chapter{Architettura}

\section{Visione d'Insieme}

Il sistema è articolato in tre macro-componenti che comunicano tra loro attraverso
interfacce ben definite:

\begin{enumerate}
  \item \textbf{Application Layer}: API REST sviluppata con FastAPI, punto di ingresso
        per tutti i client HTTP.
  \item \textbf{Data Layer}: due database specializzati — PostgreSQL per metadati
        strutturati e conversazioni, Qdrant per gli embedding vettoriali.
  \item \textbf{Intelligence Layer}: pipeline RAG orchestrata con LangGraph che
        coordina embedding, retrieval e generazione della risposta.
\end{enumerate}

\section{Diagramma a Blocchi}

\begin{figure}[H]
  \centering
  \begin{tabular}{|c|}
    \hline
    \textbf{Client (browser / curl)} \\
    \hline
  \end{tabular}
  \quad$\downarrow$ HTTP \\[0.3em]
  \begin{tabular}{|c|}
    \hline
    \textbf{FastAPI Application} \\
    \small /documents \quad /chat \quad /search \\
    \hline
  \end{tabular}

  \vspace{0.5em}
  $\swarrow$ \hspace{4cm} $\searrow$
  \vspace{0.3em}

  \begin{tabular}{c}
    \begin{tabular}{|c|}
      \hline
      \textbf{LangGraph RAG Pipeline} \\
      \small embed\_query $\to$ search\_qdrant \\
      \small build\_prompt $\to$ call\_llm \\
      \hline
    \end{tabular}
  \end{tabular}
  \hspace{1cm}
  \begin{tabular}{c}
    \begin{tabular}{|c|}
      \hline
      \textbf{PostgreSQL} \\
      \small documents, conversations \\
      \hline
    \end{tabular}
  \end{tabular}

  \vspace{0.5em}
  $\downarrow$ \hspace{0cm}
  $\downarrow$

  \begin{tabular}{|c|} \hline \textbf{Qdrant} \\ \small doc\_chunks \\ \hline \end{tabular}
  \hspace{1cm}
  \begin{tabular}{|c|} \hline \textbf{OpenAI / Ollama} \\ \small LLM + Embeddings \\ \hline \end{tabular}

  \caption{Architettura a blocchi dell'AI Document Assistant}
\end{figure}

\section{Flusso di Ingestion (Upload PDF)}

\begin{enumerate}
  \item Il client invia il file tramite \texttt{POST /documents/upload}.
  \item FastAPI valida il file (mime type, dimensione, header PDF).
  \item \texttt{ingestion.py} estrae il testo pagina per pagina con \textbf{pypdf}.
  \item Il testo viene suddiviso in chunk da 500 caratteri con overlap di 50
        mediante \texttt{RecursiveCharacterTextSplitter} di LangChain.
  \item Per ogni chunk viene generato un vettore di embedding (1536 dimensioni
        con \texttt{text-embedding-3-small}).
  \item I vettori sono inseriti in Qdrant nella collection \texttt{doc\_chunks}
        con payload: \texttt{doc\_id}, \texttt{page\_num}, \texttt{chunk\_index}, \texttt{text}.
  \item I metadati del documento (id, filename, total\_chunks, total\_pages, status)
        vengono salvati in PostgreSQL.
\end{enumerate}

\section{Flusso RAG (Chat)}

Il flusso RAG è implementato come \textbf{grafo deterministico} con LangGraph.
Il grafo ha quattro nodi connessi sequenzialmente:

\begin{table}[H]
  \centering
  \begin{tabularx}{\textwidth}{llX}
    \toprule
    \textbf{Nodo} & \textbf{Funzione} & \textbf{Operazione} \\
    \midrule
    \texttt{embed\_query}   & Embedding        & Converte la domanda in vettore \\
    \texttt{search\_qdrant} & Retrieval        & Cerca i top-k chunk simili \\
    \texttt{build\_prompt}  & Prompt Building  & Compone il prompt con contesto \\
    \texttt{call\_llm}      & Generation       & Genera la risposta finale \\
    \bottomrule
  \end{tabularx}
  \caption{Nodi del grafo LangGraph}
\end{table}

Lo stato condiviso tra i nodi è un dizionario tipizzato (\texttt{TypedDict}) con
i campi: \texttt{question}, \texttt{doc\_id}, \texttt{top\_k}, \texttt{query\_vector},
\texttt{retrieved\_chunks}, \texttt{prompt\_messages}, \texttt{answer}.

\section{Motivazione per LangGraph}

LangGraph è stato scelto per le seguenti ragioni:

\begin{itemize}
  \item \textbf{Controllo esplicito del flusso}: ogni transizione tra nodi è dichiarata,
        rendendo il pipeline auditabile e modificabile senza riscrivere la logica.
  \item \textbf{Estendibilità}: aggiungere un nodo di reranking, di verifica delle fonti
        o di tool-calling richiede solo l'aggiunta di un nodo e di un edge nel grafo.
  \item \textbf{Compatibilità}: si integra nativamente con LangChain, riutilizzando
        \texttt{ChatOpenAI}, \texttt{OllamaEmbeddings}, ecc.
  \item \textbf{Streaming}: il grafo supporta lo streaming token-by-token,
        implementato nell'endpoint \texttt{/chat} con Server-Sent Events.
\end{itemize}

\section{Modello Dati}

\subsection{PostgreSQL}

\begin{lstlisting}[language=SQL, caption=Schema tabelle PostgreSQL]
-- Metadati documenti
CREATE TABLE documents (
    id             VARCHAR(36) PRIMARY KEY,
    filename       VARCHAR(512) NOT NULL,
    original_filename VARCHAR(512) NOT NULL,
    upload_date    TIMESTAMPTZ DEFAULT NOW(),
    total_chunks   INTEGER DEFAULT 0,
    total_pages    INTEGER DEFAULT 0,
    file_size_bytes INTEGER DEFAULT 0,
    status         VARCHAR(32) DEFAULT 'processing'
);

-- Storico conversazioni
CREATE TABLE conversations (
    id         VARCHAR(36) PRIMARY KEY,
    session_id VARCHAR(36) NOT NULL,
    doc_id     VARCHAR(36) REFERENCES documents(id) ON DELETE SET NULL,
    role       VARCHAR(16) NOT NULL,  -- 'user' | 'assistant'
    content    TEXT NOT NULL,
    created_at TIMESTAMPTZ DEFAULT NOW()
);
\end{lstlisting}

\subsection{Qdrant Collection}

Ogni punto nella collection \texttt{doc\_chunks} ha la seguente struttura:

\begin{lstlisting}[language=Python, caption=Struttura punto Qdrant]
PointStruct(
    id=f"{doc_id}_{chunk_index}",   # ID stringa
    vector=[0.023, -0.415, ...],    # 1536 float32
    payload={
        "doc_id": "uuid-string",
        "chunk_index": 42,
        "page_num": 3,
        "text": "...testo del chunk..."
    }
)
\end{lstlisting}
